% Werke X/1 (1917), native184–187, printed177–180.
\begingroup
\setlength{\headheight}{25pt}
\makeatletter\tagsleft@true\let\veqno\@@leqno\makeatother
\SourceLeaf{184}\NoSourceHead
\SourceGroup{II}
\HistoricTitle{[REIHENENTWICKELUNGEN UND BEZIEHUNGEN\\ ZUM ELLIPSENUMFANG.]}
\HistoricCitation{[Eintragungen im \textsc{Leiste}.]}
\NumberedItem{1}
\noindent [S. 25]
\begin{equation}\tag*{[1]}1-\frac{1.1}{2.2}BB-\frac{1.1.1.3}{2.2.4.4}B^4-\text{ etc.}\end{equation}
\begin{equation}\tag*{[2]}B+\frac14B^3+\frac9{64}B^5\ldots=4z+8z^5\end{equation}
\begin{equation}\tag*{[3]}B=4z-16z^3-376z^5\end{equation}
\begin{equation}\tag*{[4]}\text{Periph.}=1-4z^2+20z^4+800z^6\end{equation}
\ClosingRule
\noindent [S. 26]
\begin{equation}\tag*{[5]}\begin{gathered}
\text{Peripheria Ellipseos}\\
=1-\frac{1.1}{2.2}BB-\frac{1.1.1.3}{2.2.4.4}B^4\text{ etc.}
\end{gathered}\end{equation}
\begin{equation}\tag*{[6]}B+\frac14B^3+\frac14.\frac9{16}B^5\ldots=(2z^{\frac12}+2z^{\frac92}+\text{etc.})^2=rr\end{equation}
\begin{equation}\tag*{[7]}=\frac{2bbrdr}{BBdB}-\frac{rr}{B^3}\end{equation}
\begin{equation}\tag*{[8]}\frac{2Bbbrdr-rrdB}{B^3dB}=C\frac{r^4bbxdp}{B^4pdx}\end{equation}
\[\frac{2bbrdr}{BBdB}-\frac{bb}{BB}\qquad\frac{bb}{BdB}\left(\frac{2dr}{r}-\frac{dB}{B}\right)\]
\begin{equation}\tag*{[9]}\frac{2BBbbdr-rBdB}{r^3bb[dB]}=C\frac{xdp}{pdx}.\end{equation}
\par\vfill\noindent X1.\hfill23
\clearpage
\SourceLeaf{185}\ArithmeticHead{178}{L}{ANALYSIS.\quad NACHLASS.}
\NumberedItem{2}
\noindent [S. 37]
\begin{equation}\tag*{[10]}\int\frac{\frac{bb}{a}+\left(\frac{aa-bb}{a}\right)ss}{\surd\left(1-\frac{aa-bb}{aa}ss\right)}[ds].\end{equation}
\ClosingRule
\begin{equation}\tag*{[11]}\mathbf{M}\nu=\frac12+\frac12\frac{1-\surd(1-\nu\nu)}{1+\surd(1-\nu\nu)}\mathbf{M}\frac{1-\surd(1-\nu\nu)}{1+\surd(1-\nu\nu)}\end{equation}
\begin{equation}\tag*{[12]}\mathbf{M}\frac{2\surd x}{1+x}=\frac12+\frac12x\mathbf{M}x\end{equation}
\begin{equation}\tag*{[13]}\mathbf{M}x=\frac12+\frac1{16}xx+\frac1{32}x^4+\frac{41}{2048}x^6+\ldots\end{equation}
\ClosingRule
\begin{center}Peripheria Ellipsis cuius axes \(a,b\) est\end{center}
\begin{equation}\tag*{[14]}\pi\frac{bb}{a}\left|K\surd\frac{aa-bb}{aa}\right|\left(1+\frac{aa-bb}{bb}\mathbf{M}\surd\frac{aa-bb}{aa}\right)\end{equation}
\begin{equation}\tag*{[15]}=\pi\frac{bb}{a}\frac{1+N\surd\frac{aa-bb}{bb}}{\text{Med. inter 1 et }\frac ba}\end{equation}
\begin{equation}\tag*{[16]}N\nu=\nu\nu\mathbf{M}\nu=\frac12\nu\nu+\nu.N\frac{1-\surd(1-\nu\nu)}{1+\surd(1-\nu\nu)}.\end{equation}
\NumberedItem{3}
\noindent S. 48]
\begin{equation}\tag*{[17]}q=\frac\pi2a\left(1-\frac{1.1}{2.2}\left(\frac{aa-bb}{aa}\right)-\frac{1.1.1.3}{2.2.4.4}\left(\frac{aa-bb}{aa}\right)^2\ldots\right)\end{equation}
\begin{equation}\tag*{[18]}\frac{aa}{M(a,b)}=a\left(1+\frac{1.1}{2.2}\left(\frac{aa-bb}{aa}\right)+\frac{1.1.3.3}{2.2.4.4}\left(\frac{aa-bb}{aa}\right)^2\ldots\right)\end{equation}
\begin{equation}\tag*{[19]}\frac{a}{M(a,b)}-\frac{2q}{\pi a}=\frac12\left(\frac{aa-bb}{aa}\right)+\frac{1.1.3}{2.2.4}\left(\frac{aa-bb}{aa}\right)^2\ldots\end{equation}
\ClosingRule
\clearpage
\SourceLeaf{186}\fancyhead{}\fancyhead[R]{\small179}\fancyhead[C]{\scriptsize REIHEN.\quad BEZIEHUNGEN ZUM ELLIPSENUMFANG.}
\NumberedItem{4}
\noindent [S. 49]
\begin{equation}\tag*{[20]}K\frac{2x^{\frac12}}{1+x}=(1+x)Kx\end{equation}
\begin{equation}\tag*{[21]}K\surd(1-xx)=Hx\qquad\frac{1-x}{1+x}=y\end{equation}
\begin{equation}\tag*{[22]}K\frac{1-x}{1+x}=H\frac{2x^{\frac12}}{1+x}\end{equation}
\begin{equation}\tag*{[23]}\frac{2y^{\frac12}}{1+y}=\surd(1-xx)\end{equation}
\begin{equation}\tag*{[24]}H\frac{2x^{\frac12}}{1+x}=\frac{1+x}{2}Hx\end{equation}
\begin{equation}\tag*{[25]}Hx=\frac2{1+x}H\frac{2x^{\frac12}}{1+x}\end{equation}
\begin{equation}\tag*{[26]}\frac{2x^{\frac12}}{1+x}={}'x\end{equation}
\begin{equation}\tag*{[27]}\frac{1-\surd(1-xx)}{1+\surd(1-xx)}=x'\end{equation}
\begin{equation}\tag*{[28]}\frac12+\frac14x'+\frac18x'x''+\text{etc.}=\mathbf{M}x.\end{equation}
\ClosingRule
\NumberedItem{5}
\noindent S. 83]
\begin{equation}\tag*{[29]}L\frac{2\nu^{\frac12}}{1+\nu}=\frac{1+\nu}{2}\nu L\nu+\frac{\nu+1}{2}K\nu\qquad\frac LK=\mathbf M\end{equation}
\begin{equation}\tag*{[30]}\mathbf M\frac{2\nu^{\frac12}}{1+\nu}=\frac12+\frac12\nu\mathbf M\nu\end{equation}
\begin{equation}\tag*{[31]}1)\quad\frac2{1+\nu}\left(K\frac{2\nu^{\frac12}}{1+\nu}-L\frac{2\nu^{\frac12}}{1+\nu}\right)=K\nu-\nu L\nu\end{equation}
\begin{flushright}(Verif.)\end{flushright}
\begin{equation}\tag*{[32]}2)\qquad K\frac{2\nu^{\frac12}}{1+\nu}=(1+\nu)K\nu.\end{equation}
\ClosingRule
\par\vfill\hfill23\textsuperscript{*}
\clearpage
\SourceLeaf{187}\ArithmeticHead{180}{L}{ANALYSIS.\quad NACHLASS.}
\noindent [S. 88]
\begin{equation}\tag*{[33]}1+\frac12.\frac12\nu\nu+\frac12.\frac12.\frac34.\frac34\nu^4+\text{etc.}=K\end{equation}
\begin{equation}\tag*{[34]}\frac12\nu+\frac12.\frac12.\frac34\nu^3+\frac{1.1.3.3.5}{2.2.4.4.6}\nu^5+\text{etc.}=\mathbf L\end{equation}
\begin{equation}\tag*{[35]}K\nu=\frac1{M\surd(1-\nu\nu)}=\frac1{\overline{1+\nu}M\frac{1-\nu}{1+\nu}}\end{equation}
\begin{equation}\tag*{[36]}\frac{\partial K}{\partial\nu}=\nu\frac{\partial\mathbf L}{\partial\nu},\quad\frac{\nu\partial\nu K}{\partial\nu}=\frac{\partial\nu\mathbf L}{\partial\nu}\end{equation}
\begin{equation}\tag*{[37]}K'=\nu\mathbf L',\quad\nu K+\nu\nu K'=\mathbf L+\nu\mathbf L'\end{equation}
\begin{equation}\tag*{[38]}K+\left(\nu-\frac1\nu\right)K'=L\end{equation}
\begin{equation}\tag*{[39]}(1+\nu)K\nu=K\frac{2\sqrt\nu}{1+\nu}\end{equation}
\begin{equation}\tag*{[40]}(1+\nu)K'\nu+K\nu=K'\frac{2\sqrt\nu}{1+\nu}\Big\}\times\frac{1-\nu}{(1+\nu)^2\nu^{\frac12}}.\end{equation}
\NumberedItem{6}
\noindent S. 56]
\[\renewcommand{\arraystretch}{1.2}\begin{array}{r@{\quad}l@{\quad}l@{\quad}l}
\ldots{}'a&a&a'&a''\\
{}'b&b&b'&b''\\
{}'A&A&A'&A''\\
{}'B&B&B'&B''
\end{array}\qquad\begin{array}{l@{\;}l}
A=a+b&\\
B=a-b&\\\cline{1-1}
A'=a&=\tfrac12({}'a+{}'b)\\
B'=\surd(aa-bb)&=\tfrac12({}'a-{}'b)\\\hline
A''=\tfrac14({}''a+{}''b)&\\
B''=\tfrac14({}''a-{}''b)&
\end{array}\]
\[\renewcommand{\arraystretch}{1.3}\begin{array}{l@{\qquad}l}
3{,}38889961&0{,}025313952\\
1{,}70710678&0{,}29289322\\
1&0{,}70710675\ [=\surd\tfrac12]\\
0{,}8535533912&0{,}8408964152
\end{array}\]
\ClosingRule
\clearpage
\endgroup
